\documentclass[12pt,letterpaper]{exam}
\usepackage[lmargin=1in,rmargin=1in,tmargin=1in,bmargin=1in]{geometry}
\usepackage{../style/exams}

% -------------------
% Course & Exam Information
% -------------------
\newcommand{\course}{MATH 344: Final Exam}
\newcommand{\term}{Spring ---\textsubscript{\textsubscript{1}} 2026}
\newcommand{\examdate}{04/29/2026}
\newcommand{\timelimit}{150 Minutes}

\setbool{hideans}{true} % Student: True; Instructor: False

\usepackage{array}
\newcolumntype{C}[1]{>{\centering\arraybackslash}p{#1}}
\newcommand{\clinevsol}[1]{\ifbool{hideans}{}{#1}}
\newcommand{\vanbox}[1]{\ifbool{hideans}{#1}{\boxed{#1}}}
\DeclareMathOperator{\rank}{rank}
\DeclareMathOperator{\nullity}{nullity}
\DeclareMathOperator{\mspan}{span}

\usepackage{arydshln} % Dotted Augmented Matrix - Load After 'array' if used

% -------------------
% Content
% -------------------
\begin{document}

\examtitle
\instructions{Write your name on the appropriate line on the exam cover sheet. This exam contains \numpages\ pages (including this cover page) and \numquestions\ questions. Check that you have every page of the exam. Answer the questions in the spaces provided on the question sheets. Be sure to answer every part of each question and show all your work. If you run out of room for an answer, continue on the back of the page --- being sure to indicate the problem number.} 
\scores
\bottomline
\newpage


% -------------------
% Questions
% -------------------
\begin{questions}

% ------------- Exam 1 ------------- %

% Question 1
\newpage
\question Let $\mathbf{u}= \langle \sqrt{3}, 1 \rangle$ and $\mathbf{v}= \langle \sqrt{3}, 3 \rangle$. \par\vspace{0.2cm}
	\begin{parts}
	% (a)
	\part[2] Compute $\mathbf{u} - 2\mathbf{v}$. \par\vspace{0.3cm}
		\[
		\psol{\mathbf{u} - 2\mathbf{v}= \langle \sqrt{3}, 1 \rangle - 2 \langle \sqrt{3}, 3 \rangle = \langle \sqrt{3}, 1 \rangle - \langle 2 \sqrt{3}, 6 \rangle = \langle \sqrt{3} - 2\sqrt{3}, 1 - 6 \rangle= \boxed{\langle -\sqrt{3}, -5 \rangle}}
		\] \par\vspace{0.5cm}\vfill
	
	% (b)
	\part[2] Are $\mathbf{u}$ and $\mathbf{v}$ perpendicular? Justify your answer. \par\vspace{1.2cm}
	\psol{\itshape We know that two vectors are perpendicular if and only if their dot product is $0$. We have\dots}
		\[
		\psol{\mathbf{u} \cdot \mathbf{v}= \langle \sqrt{3}, 1 \rangle \cdot \langle \sqrt{3}, 3 \rangle= \sqrt{3} (\sqrt{3} + 1 (3)= 3 + 3= 6 \neq 0}
		\]
	\psol{\itshape Therefore, $\mathbf{u}$ is not perpendicular to $\mathbf{v}$.} \par\vspace{1cm}\vfill
	
	% (c)
	\part[5] Find the angle between $\mathbf{u}$ and $\mathbf{v}$. \par\vspace{0.2cm}
	\psol{\itshape We know that $\mathbf{u} \cdot \mathbf{v}= \|\mathbf{u} \| \, \|\mathbf{v}\| \, \cos\theta$. We know $\|\mathbf{u}\|= \sqrt{(\sqrt{3})^2 + 1^2}= \sqrt{3 + 1}= \sqrt{4}= 2$ and $\|\mathbf{v}\|= \sqrt{(\sqrt{3})^2 + 3^2}= \sqrt{3 + 9}= \sqrt{12}= \sqrt{4 \cdot 3}= 2 \sqrt{3}$. From (b), we know $\mathbf{u} \cdot \mathbf{v}= 6$. But then\dots}
		\[
		\begin{gathered}
		\psol{\mathbf{u} \cdot \mathbf{v}= \|\mathbf{u} \| \, \|\mathbf{v}\| \, \cos\theta} \\[0.1cm]
		\psol{6= 2 \cdot 2\sqrt{3} \cos \theta} \\[0.1cm]
		\psol{6= 4 \sqrt{3} \cos \theta} \\[0.1cm]
		\psol{\cos \theta= \tfrac{6}{4 \sqrt{3}}} \\[0.1cm]
		\psol{\cos \theta= \tfrac{\sqrt{3}}{2}} \\[0.1cm]
		\psol{\theta= \arccos \left( \tfrac{\sqrt{3}}{2} \right)} \\[0.1cm]
		\psol{\boxed{\theta= \tfrac{\pi}{6}= 30^\circ}}
		\end{gathered}
		\] \vfill
	
	% (d)
	\part[1] Find a unit vector parallel to $\mathbf{u}$. \par\vspace{0.2cm}
		\[
		\psol{\dfrac{\mathbf{u}}{\|\mathbf{u}\|}= \dfrac{\langle \sqrt{3}, 1 \rangle}{2}= \boxed{\left\langle \dfrac{\sqrt{3}}{2}, \dfrac{1}{2} \right\rangle}}
		\] \vfill
	\end{parts}



% Question 2
\newpage
\question Consider the matrices $A, B$ and the vector $\mathbf{u}$ given below.
	\[
	A= \begin{pmatrix} 1 & 0 & 2 \\ -1 & 2 & 0 \\ 1 & 3 & 1 \end{pmatrix}, \quad B= \begin{pmatrix} -1 & 0 & -1 & 6 & 4 \\ 3 & 1 & 0 & 1 & 2 \\ 2 & 6 & 1 & -3 & 0 \end{pmatrix}, \quad \mathbf{u}= \begin{pmatrix} -2 \\ 2 \\ 1 \end{pmatrix}
	\] \par\vspace{0.2cm}
	
	\begin{parts}
	% (a)
	\part[3] Compute the $(2,4)$-entry of $AB$, $(AB)_{24}$, \textit{without} explicitly computing $AB$ entirely. \par\vspace{0.2cm}
	\psol{\itshape We know the $(i,j)$-entry of a product of matrices is the $i$th row of the first dotted with the $j$th column of the second. Therefore, we have\dots}
		\[
		\psol{\langle -1, 2, 0 \rangle \cdot \langle 6, 1, -3 \rangle= -1(6) + 2(1) + 0(-3)= -6 + 2 + 0= \boxed{-4}}
		\] \vfill
	
	% (b)
	\part[3] Compute $A\mathbf{u}$ by treating the product as a linear combination of columns of A. \par\vspace{0.2cm}
		\[
		\psol{\hspace{-1.5cm} A\mathbf{u}= \begin{pmatrix} 1 & 0 & 2 \\ -1 & 2 & 0 \\ 1 & 3 & 1 \end{pmatrix} \begin{pmatrix} -2 \\ 2 \\ 1 \end{pmatrix}= -2 \begin{pmatrix} 1 \\ -1 \\ 1 \end{pmatrix} + 2 \begin{pmatrix} 0 \\ 2 \\ 3 \end{pmatrix} + 1 \begin{pmatrix} 2 \\ 0 \\ 1 \end{pmatrix}= \begin{pmatrix} -2 \\ 2 \\ -2 \end{pmatrix} + \begin{pmatrix} 0 \\ 4 \\ 6 \end{pmatrix} + \begin{pmatrix} 2 \\ 0 \\ 1 \end{pmatrix}= \boxed{\begin{pmatrix} 0 \\ 6 \\ 5 \end{pmatrix}}}
		\] \par\vspace{1cm} \vfill
	
	% (c)
	\part[4] Compute $\det A$. \par\vspace{0.2cm}
	\psol{\itshape Using cofactor expansion along the first row, we have\dots}
		\[
		\begin{gathered}
		\psol{\det A= 1 \begin{vmatrix} 2 & 0 \\ 3 & 1 \end{vmatrix} - 0 \begin{vmatrix} -1 & 0 \\ 1 & 1 \end{vmatrix} + 2 \begin{vmatrix} -1 & 2 \\ 1 & 3 \end{vmatrix}} \\[0.2cm]
		\psol{\det A= 1 \big( 2(1) - 3(0) \big) - 0 + 2 \big( -1(3) - 1(2) \big)} \\[0.2cm]
		\psol{\det A= 1 \big(2 - 0 \big) - 0 + 2 \big( -3 - 2 \big)} \\[0.2cm]
		\psol{\det A= 1(2) - 0 + 2(-5)} \\[0.2cm]
		\psol{\det A= 2 - 0 - 10} \\[0.2cm]
		\psol{\boxed{\det A= -8}}
		\end{gathered}
		\] \par\vspace{1cm} \vfill
	\end{parts}



% Question 3
\newpage
\question Suppose matrices $A, B$ both had the following row operations performed on it: \par
	\begin{table}[!ht]
	\centering
	\begin{tabular}{c}
	$5R_1 \to R_1$ \\[0.1cm]
	$R_2 \longleftrightarrow R_3$ \\[0.1cm]
	$R_1 + R_2 \to R_2$ \\[0.1cm]	
	$R_1 - R_3 \to R_3$ \\[0.1cm]
	$2R_2 + R_3 \to R_3$
	\end{tabular}
	\end{table} \par
After these row operations, the matrices $A', B'$ below, respectively, were obtained:
	\[
	A'= \begin{pmatrix} 10 & -2 & 4 \\ 0 & -1 & 6 \\ 0 & 0 & 3 \end{pmatrix} \qquad B' = \begin{pmatrix} \vanbox{-3} & 0 & 1 & 4 & 0 \\ 0 & \vanbox{1} & 5 & -1 & 0 \\ 0 & 0 & 0 & 0 & 0 \end{pmatrix}
	\]

\begin{parts}
% (a)
\part[3] Compute $\det A$. \par\vspace{0.2cm}
\psol{\itshape The matrix $A'$ is triangular, so its determinant is the product of the diagonal entries. We can keep track of the effects of the row operations on $A$. Undoing those effects on the determinant of $A'$ (which we include ``on the left'' below in the order of the row operations), we obtain $\det A$:}
	\[
	\psol{\det A= \left( \frac{1}{5} \cdot -1 \cdot 1 \cdot -1 \cdot 1 \right) \cdot (10 \cdot -1 \cdot 3)= \frac{1}{5} \cdot -30= \boxed{-6}}
	\] \vfill

% (b)
\part[2] Using (a), determine whether $A$ is invertible, i.e. whether $A^{-1}$ exists. \par\vspace{0.2cm}
\psol{\itshape We know a matrix is invertible if and only if its determinant is not $0$. We know from (a) that $\det A= -6 \neq 0$. Therefore, $A$ is invertible, i.e. $A^{-1}$ exists.} \par\vspace{1cm}\vfill

% (c)
\part[2] Determine the rank and nullity of $B$. \par\vspace{0.2cm}
\psol{\itshape Observe that $B'$ is in row-echelon form (though not RREF). As $B'$ was obtained from $B$ using elementary row operations, $B$ and $B'$ have the same rank and nullity. We know the rank is the number of pivot positions. We see $\rank B'= 2$ from above. The number of non-pivot columns is the nullity. So, $\nullity B'= 3$. Therefore, we have\dots}
	\[
	\psol{\boxed{\rank B= 2 \text{ and } \nullity B= 3}}
	\] \vfill

% (d)
\part[3] Does $B\mathbf{x}= \mathbf{0}$ have a solution? If so, is the solution unique? \par\vspace{0.2cm}
\psol{\itshape If $A \mathbf{x}= \mathbf{b}$ has a solution, we know the solution is unique if and only if $\nullity A= 0$. The matrix-vector equation $A\mathbf{x}= \mathbf{0}$ always has at least one solution---the zero vector. We know $\nullity B= 3$. Therefore, $B\mathbf{x}= \mathbf{0}$ has infinitely many solutions.} \par\vspace{1cm}
\end{parts}



% Question 4
\newpage
\question[10] Matrices $M, N, P$ below are in RREF and represent augmented matrices coming from systems of linear equations. For each matrix, determine whether the corresponding system is consistent or inconsistent. If it is inconsistent, explain why. If it is consistent, either find the unique solution or parametrize the solution set (if infinitely many solutions exist). Use variables $x_1, x_2, \ldots$\ .
	\[
	M= \begin{pmatrix} 1 & -1 & 0 \\ 0 & 0 & 0 \end{pmatrix}, \quad N= \begin{pmatrix} 1& 0 & 0 \\ 0 & 0 & 1 \end{pmatrix}, \quad P= \begin{pmatrix} 1 & 0 & 0 \\ 0 & 1 & 0 \end{pmatrix}
	\] \par\vspace{0.2cm}

\begin{enumerate}[(a)]
% M
\item System Corresponding to $M$: \psol{\itshape Including a vertical line to separate coefficients from `solutions', we have\dots}
	\[
	\psol{%
	\left[
	\begin{array}{cc|c}
	\vanbox{1} & -1 & 0 \\
	0 & 0 & 0 \\[0.1pt]
	\end{array}
	\right]}
	\]
\psol{\itshape Labeling our variables as $x_1$ and $x_2$, we can see that $x_1$ is a basic variable and $x_2$ is a free variable. We can see from the first row that $x_1 - x_2= 0$, i.e. $x_1= x_2$. Therefore, the solutions are\dots}
	\[
	\psol{\begin{pmatrix} x_1 \\ x_2 \end{pmatrix}= \begin{pmatrix} x_2 \\ x_2 \end{pmatrix}= \boxed{x_2 \begin{pmatrix} 1 \\ 1 \end{pmatrix}}}
	\]
\psol{\itshape where $x_2$ is any real number.} \vfill

% N
\item System Corresponding to $N$: \psol{\itshape Including a vertical line to separate coefficients from `solutions', we have\dots}
	\[
	\psol{%
	\left[
	\begin{array}{cc|c}
	1 & 0 & 0 \\
	0 & 0 & 1 \\[0.1pt]
	\end{array}
	\right]}
	\]
\psol{\itshape From the second row, we see that $0= 1$, which is impossible. Therefore, the system is inconsistent, i.e. there are no solutions.} \par\vspace{1cm} \vfill

% P 
\item System Corresponding to $P$: \psol{\itshape Including a vertical line to separate coefficients from `solutions', we have\dots}
	\[
	\psol{%
	\left[
	\begin{array}{cc|c}
	1 & 0 & 0 \\
	0 & 1 & 0 \\[0.1pt]
	\end{array}
	\right]}
	\]
\psol{\itshape Labeling our variables as $x_1$ and $x_2$, we can see from the first row that $1x_1 + 0x_2= 0$, i.e. $x_1= 0$. From the second row, we can see $0x_1 + 1x_2= 0$, i.e. $x_2= 0$. Therefore, there is a unique solution and it is $(x_1, x_2)= (0, 0)$.}
	\[
	\psol{\boxed{(x_1, x_2)= (0, 0)}}
	\] \vfill
\end{enumerate}



% Question 5
\newpage
\question[10] Write the matrix-vector equation corresponding to the system of equations below. Solve this system using the inverse of the coefficient matrix. 
	\[
	\begin{cases}
	2x + 5y= 0 \\
	x - y= -7
	\end{cases}
	\] \par\vspace{0.2cm}

\vsol{\itshape We know the corresponding matrix-vector equation is $A\mathbf{x}= \mathbf{b}$, where $A$ is the coefficient matrix, $\mathbf{x}$ is the `vector of variables', and $\mathbf{b}$ is the vector of `solutions.' Therefore, the corresponding matrix-vector equation is\dots
	\[
	\boxed{\underbrace{\begin{pmatrix} 2 & 5 \\ 1 & -1 \end{pmatrix}}_{A} \underbrace{\begin{pmatrix} x \\ y \end{pmatrix}}_{\mathbf{x}}= \underbrace{\begin{pmatrix} 0 \\ -7 \end{pmatrix}}_{\mathbf{b}}}
	\]	
Observe that\dots
	\[
	\det A= 2(-1) - 1(5)= -2 - 5= -7 \neq 0
	\]
Therefore, $A$ is invertible, i.e. $A^{-1}$ exists. If $B= \begin{pmatrix} a & b \\ c & d \end{pmatrix}$ is invertible, we know that $B^{-1}= \frac{1}{\det B} \begin{pmatrix} d & -b \\ -c & a \end{pmatrix}$. Therefore, we have\dots
	\[
	A^{-1}= \frac{1}{\det A} \begin{pmatrix} d & -b \\ -c & a \end{pmatrix}= \frac{1}{-7} \begin{pmatrix} -1 & -5 \\ -1 & 2 \end{pmatrix}
	\]
Therefore, we have\dots
	\[
	\begin{gathered}
	\begin{pmatrix} 2 & 5 \\ 1 & -1 \end{pmatrix} \begin{pmatrix} x \\ y \end{pmatrix}= \begin{pmatrix} 0 \\ -7 \end{pmatrix} \\[0.3cm]
	\frac{1}{-7} \begin{pmatrix} -1 & -5 \\ -1 & 2 \end{pmatrix} \begin{pmatrix} 2 & 5 \\ 1 & -1 \end{pmatrix} \begin{pmatrix} x \\ y \end{pmatrix}= \frac{1}{-7} \begin{pmatrix} -1 & -5 \\ -1 & 2 \end{pmatrix} \begin{pmatrix} 0 \\ -7 \end{pmatrix} \\[0.3cm]
	\begin{pmatrix} 1 & 0 \\ 0 & 1 \end{pmatrix} \begin{pmatrix} x \\ y \end{pmatrix}= \dfrac{1}{-7} \begin{pmatrix} -1(0) + (-5)(-7) \\ -1(0) + 2(-7) \end{pmatrix} \\[0.3cm]
	\begin{pmatrix} x \\ y \end{pmatrix}= \dfrac{1}{-7} \begin{pmatrix} 0 + 35 \\ 0 - 14 \end{pmatrix} \\[0.3cm]
	\begin{pmatrix} x \\ y \end{pmatrix}= \dfrac{1}{-7} \begin{pmatrix} 35 \\ -14 \end{pmatrix} \\[0.3cm]
	\boxed{\begin{pmatrix} x \\ y \end{pmatrix}= \begin{pmatrix} -5 \\ 2 \end{pmatrix}} \\[0.3cm]
	\end{gathered}
	\]
}



% ------------- Exam 2 ------------- %



% Question 6
\newpage
\question[10] Find the LU decomposition of the matrix given below.
	\[
	A= \begin{pmatrix} 3 & 2 & 1 \\ 9 & 5 & 8 \\ 0 & 2 & -9 \end{pmatrix}
	\] \par\vspace{0.2cm}

\vsol{\itshape We place $A$ into REF, keeping track of the elementary row operations required. These row operations form our matrix $L$, while the resulting REF serves the role of $U$. We will then have $A= LU$. We have\dots	\par
	\begin{table}[!ht]
	\centering
	\begin{tabular}{C{3cm} C{3cm} C{3cm}} 
	\psol{Row Operations} & \psol{Current $L$} & \psol{Current $U$} \\ \clinevsol{\cline{1-3}}
	\\
	\begin{minipage}{3cm} \centering $-3R_1 + R_2 \to R_2$ \end{minipage} & $\begin{pmatrix} 1 & 0 & 0 \\ ? & 1 & 0 \\ ? & ? & 1 \end{pmatrix}$ & $\begin{pmatrix} 3 & 2 & 1 \\ 9 & 5 & 8 \\ 0 & 2 & -9 \end{pmatrix}$ \\
	\\
	$2R_2 + R_3 \to R_3$ & $\begin{pmatrix} 1 & 0 & 0 \\ 3 & 1 & 0 \\ 0 & ? & 1 \end{pmatrix}$ & $\begin{pmatrix} 3 & 2 & 1 \\ 0 & -1 & 5 \\ 0 & 2 & -9 \end{pmatrix}$ \\
	\\
	--- & $\begin{pmatrix} 1 & 0 & 0 \\ 3 & 1 & 0 \\ 0 & -2 & 1 \end{pmatrix}$ & $\begin{pmatrix} 3 & 2 & 1 \\ 0 & -1 & 5 \\ 0 & 0 & 1 \end{pmatrix}$
	\end{tabular}
	\end{table} \par
Therefore, we have\dots
	\[
	A= \;\; \underbrace{\begin{pmatrix} 1 & 0 & 0 \\ 3 & 1 & 0 \\ 0 & -2 & 1 \end{pmatrix}}_{L} \;\; \underbrace{\begin{pmatrix} 3 & 2 & 1 \\ 0 & -1 & 5 \\ 0 & 0 & 1 \end{pmatrix}}_{U}
	\]
}



% Question 7
\newpage
\question[10] A matrix $A$ can be written as $PA= LU$, where $P, L, U$ are given below.
	\[
	P= \begin{pmatrix} 0 & 1 & 0 \\ 1 & 0 & 0 \\ 0 & 0 & 1 \end{pmatrix}, \quad L= \begin{pmatrix} 1 & 0 & 0 \\ -1 & 1 & 0 \\ 2 & -1 & 1 \end{pmatrix}, \quad U= \begin{pmatrix} 2 & 0 & 1 \\ 0 & 1 & -1 \\ 0 & 0 & -4 \end{pmatrix}
	\]
Use this information to solve the equation $A\mathbf{x}= \begin{pmatrix} 5 \\ 0 \\ 3 \end{pmatrix}$. \par\vspace{0.3cm}

\vsol{\itshape Observe that the matrix $P$ merely performs the permutation $R_1 \leftrightarrow R_2$. So, multiplying by $P$ on both sides, we have\dots
	\[
	A\mathbf{x}= \begin{pmatrix} 5 \\ 0 \\ 3 \end{pmatrix} \Longrightarrow A \mathbf{x}= P \begin{pmatrix} 5 \\ 0 \\ 3 \end{pmatrix} \Longrightarrow LU \mathbf{x}= \begin{pmatrix} 0 \\ 5 \\ 3 \end{pmatrix}
	\]
Write $LU \mathbf{x}= L(U \mathbf{x})$. We know that $U\mathbf{x}$ is a vector, call this vector $\mathbf{y}$. So, we have $L\mathbf{y}= \langle 0, 5, 3 \rangle^T$. Therefore, we have\dots
	\[
	\begin{pmatrix} 1 & 0 & 0 \\ -1 & 1 & 0 \\ 2 & -1 & 1 \end{pmatrix} \begin{pmatrix} y_1 \\ y_2 \\ y_3 \end{pmatrix}= \begin{pmatrix} 0 \\ 5 \\ 3 \end{pmatrix}
	\]
We solve this system using forward-substitution:
	\[
	\begin{aligned}
	1y_1 + 0y_2 + 0y_3= 0 &\quad\Longrightarrow\quad y_1= 0 \\
	-1y_1 + 1y_2 + 0y_3= 0 &\quad\Longrightarrow\quad 0 + y_2= 5 &\quad\Longrightarrow\quad y_2= 5 \\
	2y_1 - 1y_2 + 1y_3= 3 &\quad\Longrightarrow\quad 2(0) - 1(5) + y_3= 3 &\psol{\quad\Longrightarrow\quad y_3= 8}
	\end{aligned}
	\]
Because $U\mathbf{x}= \mathbf{y}$ and $\mathbf{y}= \begin{pmatrix} 0 \\ 5 \\ 8 \end{pmatrix}$, we know\dots
	\[
	\psol{\begin{pmatrix} 2 & 0 & 1 \\ 0 & 1 & -1 \\ 0 & 0 & -4 \end{pmatrix} \begin{pmatrix} x_1 \\ x_2 \\ x_3 \end{pmatrix}= \begin{pmatrix} 0 \\ 5 \\ 8 \end{pmatrix}}
	\]
We can solve this using back-substitution:
	\[
	\begin{aligned}
	0x_1 + 0x_2 - 4x_3= 8 &\quad\Longrightarrow\quad -4x_3= 8 &\quad\Longrightarrow\quad x_3= -2 \\
	0x_1 + 1x_2 - x_3= 5 &\quad\Longrightarrow\quad x_2 + 2= 5 &\quad\Longrightarrow\quad x_2= 3 \\
	2x_1 + 0x_2 + 1x_3= 3 &\quad\Longrightarrow\quad 2x_1 - 2)= 0 &\quad\Longrightarrow\quad x_1= 1
	\end{aligned}
	\]
Therefore, the solution is $\boxed{\mathbf{x}= \begin{pmatrix} 1 \\ 3 \\ -2 \end{pmatrix}}$.}



% Question 8
\newpage
\question[10] A matrix $A$ and its reduced-row echelon form, $R$, are given below.
	\[
	A= \begin{pmatrix} 4 & -8 & 7 & 4 \\ 1 & -2 & 2 & 1 \\ 3 & -6 & 5 & 3 \end{pmatrix}, \qquad  R= \begin{pmatrix} \vanbox{1} & -2 & 0 & 1 \\ 0 & 0 & \vanbox{1} & 0 \\ 0 & 0 & 0 & 0 \end{pmatrix}
	\]

\begin{enumerate}[(a)]
\item Determine which $\mathbb{R}^n$ each of the spaces below are a subspace of and find their dimension. 
	\begin{table}[!ht]
	\centering
	\resizebox{0.7\columnwidth}{!}{%
	\begin{tabular}{|C{3cm}| C{3cm}| C{3cm}| C{3cm}|} \cline{1-4}
	& & & \\
	& $\operatorname{Col}(A)$ & $\operatorname{Row}(A)$ & $\operatorname{Null}(A)$ \\ 
	& & &\\ \cline{1-4}
	& & & \\ 
	Subspace of\dots & \psol{$\mathbb{R}^3$} & \psol{$\mathbb{R}^4$} & \psol{$\mathbb{R}^4$} \\ 
	& & & \\ \cline{1-4}
	& & & \\
	Dimension & \psol{$\rank A= 2$} & \psol{$\rank A= 2$} & \psol{$\nullity A= 2$} \\ 
	& & & \\ \cline{1-4}
	\end{tabular}%
	}
	\end{table}

\item Find a basis for $\operatorname{Col}(A)$. \par\vspace{0.2cm}
\psol{\itshape We see the pivot columns are the first and third column. So, we can use the first and third column of $A$ as basis vectors for $\operatorname{Col}(A)$.}
	\[
	\psol{\boxed{\left\{ \begin{pmatrix} 4 \\ 1 \\ 3 \end{pmatrix}, \begin{pmatrix} 7 \\ 2 \\ 5 \end{pmatrix} \right\}}}
	\]

\item Find a basis for $\operatorname{Row}(A)$. \par\vspace{0.2cm}
\psol{\itshape We see the pivot rows are the first and second rows. So, we can use the pivot rows of $A$---even of $R$---as a basis for $\operatorname{Row}(A)$. The latter is a `nicer' choice of basis.}
	\[
	\psol{\boxed{\left\{ \begin{pmatrix} 1 \\ -2 \\ 0 \\ 1 \end{pmatrix}, \begin{pmatrix} 0 \\ 0 \\ 1 \\ 0 \end{pmatrix} \right\}}}
	\]

\item Find a basis for $\operatorname{Null}(A)$. \par\vspace{0.2cm}
\psol{\itshape We need to find the solutions to $A\mathbf{x}= \mathbf{0}$. We see from $R$ that the basic variables are $x_1, x_3$ and the free variables are $x_2, x_4$. The first row tells us $x_1 - 2x_2 + x_4= 0$, i.e. $x_1= 2x_2 - x_4$, and the second row gives $x_3= 0$. So, $(x_1, x_2, x_3, x_4)= (2x_2 - x_4, x_2, 0, x_4)= (2x_2, x_2, 0, 0) + (-x_4, 0, 0, x_4)= x_2 (2, 1, 0, 0) + x_4 (-1, 0, 0, 1)$. As these vectors are linearly independent, we can take this as a basis for $\operatorname{Null}(A)$.}
	\[
	\psol{\boxed{\left\{ \begin{pmatrix} 2 \\ 1 \\ 0 \\ 0 \end{pmatrix}, \begin{pmatrix} -1 \\ 0 \\ 0 \\ 1 \end{pmatrix} \right\}}}
	\]
\end{enumerate}



% Question 9
\newpage
\question[10] Mark the statements below as True (T) or False (F). \par\vspace{1cm}

\begin{enumerate}[(a)]
\item \underline{\itshape\hspace{0.62cm}\psol{F}\hspace{0.6cm}}: The set of vectors $\{ \mathbf{v}_1, \ldots, \mathbf{v}_n \}$ is linearly independent if and only if at least one of the vectors can be written as a linear combination of the other vectors. \vfill % F

\item \underline{\itshape\hspace{0.6cm}\psol{T}\hspace{0.6cm}}: The set of vectors $\{ \mathbf{v}_1, \ldots, \mathbf{v}_n \}$ is linearly independent if and only if the equation $c_1 \mathbf{v}_1 + \cdots + c_n \mathbf{v}_n= \mathbf{0}$ has only the trivial solution. \vfill % T

\item \underline{\itshape\hspace{0.6cm}\psol{T}\hspace{0.6cm}}: If $\{ \mathbf{v}_1, \ldots, \mathbf{v}_k \}$ is a set of vectors in $\mathbb{R}^n$ with $k > n$, then the set is linearly dependent. \vfill % T

\item \underline{\itshape\hspace{0.6cm}\psol{T}\hspace{0.6cm}}: The columns of \textit{any} $6 \times 7$ matrix are linearly dependent. \vfill % T

\item \underline{\itshape\hspace{0.6cm}\psol{T}\hspace{0.6cm}}: The set $\{ \mathbf{u}, \mathbf{v} \}$ is linearly dependent if and only if $\mathbf{u}$ and $\mathbf{v}$ are parallel. \vfill % T

% \item \underline{\itshape\hspace{0.65cm}\psol{F}\hspace{0.65cm}}: 
\end{enumerate}
	


% Question 10
\newpage
\question[10] For each of the following, indicate the letter of the correct response. \par\vspace{0.3cm}

\begin{enumerate}[(i)]

\item \underline{\itshape\hspace{0.6cm}\psol{a}\hspace{0.6cm}}:  If the null space of a $9 \times 5$ matrix $A$ is $3$-dimensional, which of the following is the dimension of the column space of $A$?
	\begin{enumerate}[(a)]
	\item $2$
	\item $3$
	\item $5$
	\item $9$
	\item Cannot be determined, more information is required.
	\end{enumerate} \vfill % (a) 2

\item \underline{\itshape\hspace{0.6cm}\psol{d}\hspace{0.6cm}}: If $A$ is a $6 \times 7$ matrix with rank $6$, then how many solutions does $A \mathbf{x}= \mathbf{b}$ have for $\mathbf{b} \in \mathbb{R}^7$?
	\begin{enumerate}[(a)]
	\item None
	\item One
	\item Two
	\item Infinitely Many
	\item More information required
	\end{enumerate} \vfill % (d) Infinitely Many

\item \underline{\itshape\hspace{0.6cm}\psol{c}\hspace{0.6cm}}: If $A$ is a $3 \times 4$ matrix, then what is the largest the rank of $A$ can be?
	\begin{enumerate}[(a)]
	\item $0$
	\item $1$
	\item $3$
	\item $4$
	\item Cannot be determined.
	\end{enumerate} \vfill % (c) 3

\item \underline{\itshape\hspace{0.6cm}\psol{c}\hspace{0.6cm}}: If $A$ is a $3 \times 5$ matrix, what is the smallest the nullity of $A$ can be?
	\begin{enumerate}[(a)]
	\item $0$
	\item $1$
	\item $2$
	\item $3$
	\item $5$
	\end{enumerate} \vfill % (c) 2

\item \underline{\itshape\hspace{0.6cm}\psol{e}\hspace{0.6cm}}: If $A$ is a $4 \times 4$ matrix, which of the following would imply that $A$ is invertible?
	\begin{enumerate}[(a)]
	\item $\rank A + \nullity A= 4$
	\item $\rank A= 4$
	\item $\nullity A= 0$
	\item $\nullity A > 0$
	\item (b) \textit{and} (c)
	\end{enumerate} \vfill % (e) b and c
\end{enumerate}



% ------------- Exam 3 ------------- %



% Question 11
\newpage
\question[10] Find the eigenvalues, eigenvectors, and all associated eigenspaces for the matrix $A= \begin{pmatrix} 1 & 2 \\ 4 & 3 \end{pmatrix}$. \par\vspace{0.3cm}

\vsol{\itshape\small We first find the eigenvalues by finding the roots of the characteristic polynomial, i.e. the roots of $p(\lambda)= \det(A - \lambda I_n)$. We have\dots
	\[
	\begin{aligned}
	p(\lambda)&= \det(A - \lambda I_3) \\
	&= \det \begin{pmatrix} 1 - \lambda & 2 \\ 4 & 3 - \lambda \end{pmatrix} \\
	&= (1 - \lambda)(3 - \lambda) - 8 \\
	&= (3 - \lambda - 3 \lambda + \lambda^2) - 8 \\
	&= \lambda^2 - 4 \lambda - 5 \\
	&= (\lambda - 5)(\lambda + 1)
	\end{aligned}
	\]
If $p(\lambda)= 0$, then $(\lambda - 5)(\lambda + 1)= 0$, which implies either $\lambda - 5= 0$, i.e. $\lambda= 5$, or $\lambda + 1= 0$, i.e. $\lambda= -1$. Therefore, the eigenvalues are $\lambda= -1, 5$. 
	\[
	\boxed{\lambda= -1, 5}
	\]
We find eigenvectors, $\mathbf{v}$, associated to each eigenvalue by finding the solutions to $(A - \lambda I_n) \mathbf{v}= \mathbf{0}$, i.e. the nullspace of $A - \lambda I_n$. \par\vspace{0.3cm}

$\underline{\lambda= -1} \colon$ \par
We have\dots
	\[
	A - (-1)I_3= \begin{pmatrix} 1 - (-1) & 2 \\ 4 & 3 - (-1) \end{pmatrix}= \begin{pmatrix} 2 & 2 \\ 4 & 4 \end{pmatrix}
	\]
We can row reduce this matrix. Observe performing $-2R_1 + R_2 \to R_2$ and $\frac{1}{2} R_1 \to R_1$ obtains $\begin{pmatrix} 1 & 1 \\ 0 & 0 \end{pmatrix}$. The first row tells us $x_1 + x_2= 0$, i.e. $x_1= -x_2$. So, the solutions are $(x_1, x_2)= (-x_2, x_2)= x_2 (-1, 1)$. So, we can take $\langle -1, 1 \rangle^T$ as an eigenvector, and the associated eigenspace has as a basis\dots 
	\[
	\boxed{\left\{ \begin{pmatrix} -1 \\ 1 \end{pmatrix} \right\}}
	\]

$\underline{\lambda= 5} \colon$ \par
We have\dots
	\[
	A - 5I_3= \begin{pmatrix} 1 - 5 & 2 \\ 4 & 3 - 5 \end{pmatrix}= \begin{pmatrix} -4 & 2 \\ 4 & -2 \end{pmatrix}
	\]
We can row reduce this matrix. Observe performing $R_1 + R_2 \to R_1$ and $-\frac{1}{2} R_1 \to R_1$ obtains $\begin{pmatrix} 2 & -1 \\ 0 & 0 \end{pmatrix}$. The first row tells us $2x_1 - x_2= 0$, i.e. $x_1= \frac{1}{2} x_2$. So, the solutions are $(x_1, x_2)= (\frac{1}{2} x_2, x_2)= x_2 (\frac{1}{2}, 1)$. So, we can take $\langle \frac{1}{2}, 1 \rangle^T$ as an eigenvector. Scaling this obtains another possible eigenvector. So, scaling by $2$ yields the `nicer' eigenvector $\langle 1, 2 \rangle^T$. The associated eigenspace has as a basis\dots 
	\[
	\boxed{\left\{ \begin{pmatrix} 1 \\ 2 \end{pmatrix} \right\}}
	\]%
}



% Question 12
\newpage
\question[10] Let $A, B, C$ be matrices. The characteristic polynomials, eigenvalues, and a basis for the associated eigenspaces for each of $A, B, C$ are given below. Only one of $A, B, C$ is diagonalizable. Identify which matrix is diagonalizable and find matrices $P, D$ such that the given matrix is equal to $PDP^{-1}$. 
	\begin{table}[h!]
	\centering
	\begin{tabular}{cccc} \noalign{\vskip 2pt} \hline \noalign{\vskip 6pt}
	% Column Titles
	\textbf{Matrix} & \textbf{Characteristic Polynomial} & \textbf{Eigenvalue} & \textbf{Eigenspace Basis} \\ \noalign{\vskip 6pt} \hline \noalign{\vskip 6pt}
	% Row A
	$A$ & $p(\lambda)= (\lambda - 1)^2$ & $1$ & $\left\{ \renewcommand{\arraystretch}{1.2} \begin{pmatrix} 1 \\ 1 \end{pmatrix} \right\}$ \\ \noalign{\vskip 6pt} \hline \noalign{\vskip 6pt}
	% Row B
	$B$ & $p(\lambda)= -(\lambda + 2)^2(\lambda - 1)$ & $-2$ & $\left\{ \renewcommand{\arraystretch}{1.2} \begin{pmatrix} 1 \\ 0 \\ 2 \end{pmatrix}, \begin{pmatrix} 1 \\ 8 \\ -1 \end{pmatrix} \right\}$ \\ &  & $1$ & $\left\{ \renewcommand{\arraystretch}{1.2} \begin{pmatrix} 1 \\ -3 \\ 3 \end{pmatrix} \right\}$ \\ \noalign{\vskip 6pt} \hline \noalign{\vskip 6pt}
	% Row C
	$C$ & $p(\lambda)= (\lambda - 1)^2(\lambda - 5)^2$ & $1$ & $\left\{ \renewcommand{\arraystretch}{1.2} \begin{pmatrix} 0 \\ 1 \\ 0 \\ 0 \end{pmatrix}, \begin{pmatrix} 1 \\ 0 \\ 0 \\ 0 \end{pmatrix} \right\}$ \\ &  & $5$ & $\left\{ \renewcommand{\arraystretch}{1.2} \begin{pmatrix} 0 \\ 0 \\ 1 \\ 0 \end{pmatrix} \right\}$ \\ \noalign{\vskip 6pt} \hline \noalign{\vskip 2pt}
	\end{tabular}
	\end{table}

\vsol{\itshape\scriptsize We know a matrix $M$ is diagonalizable if and only if $M$ is similar to a diagonal matrix, i.e. $M= PDP^{-1}$ for some matrix $P$ and diagonal matrix $D$. We know this occurs if and only if $P$ is a matrix whose columns are linearly independent eigenvectors of $M$ and $D$ is a diagonal matrix consisting of the `corresponding' eigenvalues. We also know a matrix is diagonalizable if and only if the algebraic multiplicity of each eigenvalue $\lambda$ is equal to its geometric multiplicity (the dimension of the eigenspace). Equivalently, one needs to be able to find as many linearly independent eigenvectors as $M$ has columns. Note if the characteristic polynomial $p(\lambda)$ has degree $n$, then $M$ is $n \times n$. We see $A$ is $2 \times 2$ with eigenvalue $1$ with algebraic multiplicity 2 and geometric multiplicity 1, so $A$ is not diagonalizable. We see that $C$ is $4 \times 4$ with eigenvalue 1 with algebraic multiplicity 2 and geometric multiplicity 2, and eigenvalue 5 with algebraic multiplicity 2 and geometric multiplicity 1, so $C$ is not diagonalizable. Finally, $B$ is $3 \times 3$ with eigenvalue $-2$ with algebraic multiplicity 2 and geometric multiplicity 2, and eigenvalue 1 with geometric multiplicity 1, so $B$ is diagonalizable. One such diagonalization is\dots\small
	\[
	\boxed{P= \begin{pmatrix} 1 & 1 & 1 \\ 0 & 8 & -3 \\ 2 & -1 & 3 \end{pmatrix}, \quad D= \begin{pmatrix} -2 & 0 & 0 \\ 0 & -2 & 0 \\ 0 & 0 & 1 \end{pmatrix}}
	\]
\vfill
\tiny Note. One can find that $P^{-1}= \begin{pmatrix} -21 & 4 & 11 \\ 6 & -1 & -3 \\ 16 & -3 & -8 \end{pmatrix}$, so that $A= PDP^{-1}= \begin{pmatrix} 1 & 1 & 1 \\ 0 & 8 & -3 \\ 2 & -1 & 3 \end{pmatrix} \begin{pmatrix} -2 & 0 & 0 \\ 0 & -2 & 0 \\ 0 & 0 & 1 \end{pmatrix} \begin{pmatrix} -21 & 4 & 11 \\ 6 & -1 & -3 \\ 16 & -3 & -8 \end{pmatrix}= \begin{pmatrix} 42 & -8 & 11 \\ 0 & 16 & 9 \\ -64 & -6 & -24 \end{pmatrix}$.
}



% Question 13
\newpage
\question[10] Find an orthonormal basis for the subspace of $\mathbb{R}^3$ generated by the span of the vectors below.
	\[
	\left\{ \;\; \begin{pmatrix} 1 \\ 0 \\ 0 \end{pmatrix}, \;\; \begin{pmatrix} 1 \\ 0 \\ 1 \end{pmatrix}, \;\; \begin{pmatrix} 1 \\ -1 \\ 1 \end{pmatrix} \;\; \right\}
	\] \par\vspace{0.3cm}

\vsol{\itshape\small Label the vectors in the given basis (in order) as $\mathbf{u}_1, \mathbf{u}_2, \mathbf{u}_3$, respectively. We use Gram-Schmidt orthogonalization to produce an orthogonal basis $\{ \mathbf{w}_1, \mathbf{w}_2, \mathbf{w}_3 \}$. We have\dots
	\[
	\begin{aligned}
	\mathbf{v}_1&= \boxed{\mathbf{u}_1= \begin{pmatrix} 1 \\ 0 \\ 0 \end{pmatrix}} \\
	\\
	\operatorname{proj}_{\mathbf{v}_1} \mathbf{u}_2&= \dfrac{\mathbf{u}_2 \cdot \mathbf{v}_1}{\mathbf{v}_1 \cdot \mathbf{v}_1} \,\mathbf{v}_1= \dfrac{1(1) + 0(0) + 1(0)}{1(1) + 0(0) + 0(0)} \, \begin{pmatrix} 1 \\ 0 \\ 0 \end{pmatrix}= \dfrac{1}{1} \,\begin{pmatrix} 1 \\ 0 \\ 0 \end{pmatrix}= \begin{pmatrix} 1 \\ 0 \\ 0 \end{pmatrix} \\
	\mathbf{v}_2&= \mathbf{u}_2 - \operatorname{proj}_{\mathbf{v}_1} \mathbf{u}_2= \begin{pmatrix} 1 \\ 0 \\ 1 \end{pmatrix} - \begin{pmatrix} 1 \\ 0 \\ 0 \end{pmatrix}= \boxed{\begin{pmatrix} 0 \\ 0 \\ 1 \end{pmatrix}} \\
	\\
	%
	\operatorname{proj}_{\mathbf{v}_1} \mathbf{u}_3&= \dfrac{\mathbf{u}_3 \cdot \mathbf{v}_1}{\mathbf{v}_1 \cdot \mathbf{v}_1} \,\mathbf{v}_1= \dfrac{1(1) + 0(-1) + 0(1)}{1(1) + 0(0) + 0(0)} \, \begin{pmatrix} 1 \\ 0 \\ 0 \end{pmatrix}= \dfrac{1 + 0 + 0}{1 + 0 + 0}\, \begin{pmatrix} 1 \\ 0 \\ 0 \end{pmatrix}= \begin{pmatrix} 1 \\ 0 \\ 0 \end{pmatrix} \\
	\operatorname{proj}_{\mathbf{v}_2} \mathbf{u}_3&= \dfrac{\mathbf{u}_3 \cdot \mathbf{v}_2}{\mathbf{v}_2 \cdot \mathbf{v}_2} \,\mathbf{v}_2= \dfrac{1(0) + (-1)(0) + 1(1)}{0(0) + 0(0) + 1(1)} \, \begin{pmatrix} 0 \\ 0 \\ 1 \end{pmatrix}= \dfrac{0 + 0 + 1}{0 + 0 + 1} \, \begin{pmatrix} 0 \\ 0 \\ 1 \end{pmatrix} = \begin{pmatrix} 0 \\ 0 \\ 1 \end{pmatrix} \\
	\mathbf{v}_3&= \mathbf{u}_3 - \operatorname{proj}_{\mathbf{v}_1} \mathbf{u}_3 - \operatorname{proj}_{\mathbf{v}_2} \mathbf{u}_3= \begin{pmatrix} 1 \\ -1 \\ 1 \end{pmatrix} - \begin{pmatrix} 1 \\ 0 \\ 0 \end{pmatrix} - \begin{pmatrix} 0 \\ 0 \\ 1 \end{pmatrix}= \boxed{\begin{pmatrix} 0 \\ -1 \\ 0 \end{pmatrix}}
	\end{aligned}
	\]
So, $\{ \mathbf{v}_1, \mathbf{v}_2, \mathbf{v}_3 \}$ is an orthogonal basis for the subspace. However, we want an orthonormal basis, i.e. an orthogonal basis where each basis vector is a unit vector. We now need only normalize the basis vectors, i.e. take $\mathbf{w}_1= \frac{\mathbf{v}_1}{\|\mathbf{v}_1\|}$, $\mathbf{w}_2= \frac{\mathbf{v}_2}{\|\mathbf{v}_2\|}$, and $\mathbf{w}_3= \frac{\mathbf{v}_3}{\|\mathbf{v}_3\|}$. Observe\dots
	\[
	\begin{aligned}
	\| \mathbf{v}_1 \|&= \sqrt{1^2 + 0^2 + 0^2}= \sqrt{1 + 0 + 0}= \sqrt{1}= 1 \\
	\| \mathbf{v}_2 \|&= \sqrt{0^2 + 0^2 + 1^2}= \sqrt{0 + 0 + 1}= \sqrt{1}= 1 \\
	\| \mathbf{v}_3 \|&= \sqrt{0^2 + (-1)^2 + 0^2}= \sqrt{0 + 1 + 0}= \sqrt{1}= 1
	\end{aligned}
	\]
Therefore, an orthogonal basis for the subspace is\dots
	\[
	\boxed{\left\{ \begin{pmatrix} 1 \\ 0 \\ 0 \end{pmatrix}, \begin{pmatrix} 0 \\ 0 \\ 1 \end{pmatrix}, \begin{pmatrix} 0 \\ -1 \\[0.15cm] 1 \end{pmatrix} \right\}}
	\]
\tiny Note. Let the given subspace be $S$. It is clear that $\dim S= 3$, so $S= \mathbb{R}^3$. But then standard basis for $\mathbb{R}^3$ (which is orthonormal) could serve as an orthonormal basis for $S$. So, we could have chosen orthonormal basis $\left\{ \langle 1, 0, 0 \rangle^T, \langle 0, 1, 0 \rangle^T, \langle 0, 0, 1 \rangle^T \right\}$---avoiding Gram-Schmidt.}



% Question 14
\newpage
\question[10] Find the $QR$ factorization for the matrix $A= \begin{pmatrix} 3 & 1 \\ 4 & 1 \end{pmatrix}$.\pspace

\vsol{\itshape Observe the columns of $A$ are linearly independent, so $\rank A= 2$. [There are only two columns and neither is a multiple of the other, so they are not parallel and thus linearly independent.] We need to find an orthonormal basis for the column space of $A$, $\operatorname{Col} A$. We can use Gram-Schmidt to produce such a basis. Label the columns of $A$ (in order) as $\mathbf{u}_1, \mathbf{u}_2$. We have\dots
	\[
	\begin{aligned}
	\mathbf{v}_1&= \mathbf{u}_1= \begin{pmatrix} 3 \\ 4 \end{pmatrix} \\
	\\
	\operatorname{proj}_{\mathbf{v}_1} \mathbf{u}_2&= \dfrac{\mathbf{u}_2 \cdot \mathbf{v}_1}{\mathbf{v}_1 \cdot \mathbf{v}_1} \,\mathbf{v}_1= \dfrac{1(3) + 1(4)}{3(3) + 4(4)} \, \begin{pmatrix} 3 \\ 4 \end{pmatrix}= \dfrac{3 + 4}{9 + 16} \, \begin{pmatrix} 3 \\ 4 \end{pmatrix}= \dfrac{7}{25} \begin{pmatrix} 3 \\ 4 \end{pmatrix} \\
	\mathbf{v}_2&= \mathbf{u}_2 - \operatorname{proj}_{\mathbf{v}_1} \mathbf{u}_2= \begin{pmatrix} 1 \\ 1 \end{pmatrix} - \dfrac{7}{25} \begin{pmatrix} 3 \\ 4 \end{pmatrix}= \begin{pmatrix} 1 \\ 1 \end{pmatrix} - \begin{pmatrix} \tfrac{21}{25} \\[0.1cm] \frac{28}{25} \end{pmatrix}= \begin{pmatrix} \tfrac{4}{25} \\[0.1cm] -\tfrac{3}{25} \end{pmatrix}= \dfrac{1}{25} \begin{pmatrix} 4 \\ -3 \end{pmatrix}
	\end{aligned}
	\]
So, an orthogonal basis for the column space could be $\left\{ \begin{pmatrix} 3 \\ 4 \end{pmatrix}, \dfrac{1}{25} \begin{pmatrix} 4 \\ -3 \end{pmatrix} \right\}$. Of course, scaling a basis vector by a nonzero scalar still results in a basis vector. So, we scale the second basis vector by 25 to clear fractions. Therefore, we have orthogonal basis $\left\{ \begin{pmatrix} 3 \\ 4 \end{pmatrix}, \begin{pmatrix} 4 \\ -3 \end{pmatrix} \right\}$. We need now only normalize the vectors to obtain an orthonormal basis. 
	\[
	\begin{aligned}
	\| \langle 3, 4 \rangle^T \|&= \sqrt{3^2 + 4^2}= \sqrt{9 + 16}= \sqrt{25}= 5 \\
	\| \langle 4, -3 \rangle^T \|&= \sqrt{4^2 + (-3)^2}= \sqrt{16 + 9}= \sqrt{25}= 5
	\end{aligned}
	\]
So, an orthonormal basis for the column space of $A$ is $\left\{ \begin{pmatrix} \tfrac{3}{5} \\[0.1cm] \tfrac{4}{5} \end{pmatrix}, \begin{pmatrix} \tfrac{4}{5} \\[0.1cm] -\tfrac{3}{5} \end{pmatrix} \right\}$. So, we can take $Q= \begin{pmatrix} \tfrac{3}{5} & \tfrac{4}{5} \\[0.1cm] \tfrac{4}{5} & -\tfrac{3}{5} \end{pmatrix}$. Now if $A= QR$ and $Q$ is an orthogonal matrix, we know $Q^TA= Q^TQR$ and $Q^TQ= I$, so $R= Q^TA$. Thus,\dots
	\[
	R= Q^TA= \begin{pmatrix} \tfrac{3}{5} & \tfrac{4}{5} \\[0.1cm] \tfrac{4}{5} & -\tfrac{3}{5} \end{pmatrix}^T \begin{pmatrix} 3 & 1 \\ 4 & 1 \end{pmatrix}= \begin{pmatrix} \tfrac{3}{5} & \tfrac{4}{5} \\[0.1cm] \tfrac{4}{5} & -\tfrac{3}{5} \end{pmatrix} \begin{pmatrix} 3 & 1 \\ 4 & 1 \end{pmatrix}= \begin{pmatrix} \tfrac{9}{5} + \tfrac{16}{5} & \tfrac{3}{5} + \tfrac{4}{5} \\[0.1cm] \tfrac{12}{5} - \tfrac{12}{5} & \tfrac{4}{5} - \tfrac{3}{5} \end{pmatrix}= \begin{pmatrix} 5 & \tfrac{7}{5} \\[0.1cm] 0 & \tfrac{1}{5} \end{pmatrix}
	\]
Therefore, we have\dots
	\[
	\boxed{A= \underbrace{\begin{pmatrix} \tfrac{3}{5} & \tfrac{4}{5} \\[0.1cm] \tfrac{4}{5} & -\tfrac{3}{5} \end{pmatrix}}_{Q} \;\; \underbrace{\begin{pmatrix} 5 & \tfrac{7}{5} \\[0.1cm] 0 & \tfrac{1}{5} \end{pmatrix}}_{R}}
	\]
}



% Question 15
\newpage
\question[10] Find the line of best fit (the least square regression line) for the dataset given below. \par
	\begin{table}[!ht]
	\centering
	\begin{tabular}{r||ccc}
	$x$ & $-1$ & $0$ & $1$ \\ \hline
	$y$ & $-1$ & $1$ & $1$ \\[0.2pt]
	\end{tabular}
	\end{table} \par\vspace{0.3cm}

\vsol{\itshape We want to find the line of best fit (in the sense of least squares): $y= a_0 + a_1 x$. That is, we want the parameters $a_0, a_1$ such that $A \mathbf{x}= \mathbf{y}$ (which is inconsistent as the points clearly do not lie along a line) is as `close' to being consistent as possible (in the sense of least squares), where $A$ is the `matrix of model outputs', $\mathbf{x}= \langle a_0, a_1 \rangle^T$, and $\mathbf{y}$ is the vector of $y$-values. Using model $y= a_0 + a_1x$ and given data, we have\dots
	\[
	\begin{aligned}
	x= -1 &\colon a_0 + a_1(-1)= 1a_0 + (-1) a_1 \\[0.2cm]
	x= 0 &\colon a_0 + a_1(0)= 1a_0 + 0 a_1 \\[0.2cm]
	x= 1 &\colon a_0 + a_1(1)= 1a_0 + 1 a_1
	\end{aligned}
	\]
So, we have $A= \begin{pmatrix} 1 & -1 \\ 1 & 0 \\ 1 & 1 \end{pmatrix}$. We know the least squares solution is the solution to the normal equations: $A^TA \mathbf{x}= A^T \mathbf{y}$. We have\dots
	\[
	\begin{aligned}
	A^TA&= \begin{pmatrix} 1 & 1 & 1 \\ -1 & 0 & 1 \end{pmatrix} \begin{pmatrix} 1 & -1 \\ 1 & 0 \\ 1 & 1 \end{pmatrix}= \begin{pmatrix} 1 + 1 + 1 & -1 + 0 + 1 \\ -1 + 0 + 1 & 1 + 0 + 1 \end{pmatrix}= \begin{pmatrix} 3 & 0 \\ 0 & 2 \end{pmatrix} \\[0.2cm]
	A^T \mathbf{y}&= \begin{pmatrix} 1 & 1 & 1 \\ -1 & 0 & 1 \end{pmatrix} \begin{pmatrix} -1 \\ 1 \\ 1 \end{pmatrix}= \begin{pmatrix} -1 + 1 + 1 \\ 1 + 0 + 1 \end{pmatrix}= \begin{pmatrix} 1 \\ 2 \end{pmatrix}
	\end{aligned}
	\]
We now solve the normal equations:
	\[
	\begin{gathered}
	A^TA \mathbf{x}= A^T \mathbf{y} \\[0.2cm]
	\begin{pmatrix} 3 & 0 \\ 0 & 2 \end{pmatrix} \begin{pmatrix} a_0 \\ a_1 \end{pmatrix}= \begin{pmatrix} 1 \\ 2 \end{pmatrix} \\[0.2cm]
	\begin{pmatrix} 3a_0 \\ 2a_1 \end{pmatrix}= \begin{pmatrix} 1 \\ 2 \end{pmatrix}
	\end{gathered}
	\]
So, $3a_0= 1$, i.e. $a_0= \tfrac{1}{3}$, and $2a_1= 2$, i.e. $a_1= 1$. So, our model is $y= a_0 + a_1x= \tfrac{1}{3} + 1x$. Therefore, the line of best fit is\dots
	\[
	\boxed{y= x + \dfrac{1}{3}}
	\] \vfill

\tiny Note. One can show that the error vector is $\mathbf{y} - \mathbf{y}_{\text{pred}}= \langle -\tfrac{1}{3}, \tfrac{2}{3}, -\tfrac{1}{3} \rangle^T$, so the least squares error is $\| \mathbf{y} - \mathbf{y}_{\text{pred}} \|= \sqrt{\tfrac{2}{3}} \approx 0.816497$.}

\end{questions}
\end{document}